\section{综合框架与关键结论}

数据并行 CUDA 程序的基本结构可以归纳为以下完整链条：
\begin{enumerate}
  \item 从数据依赖出发，定义每个线程的独立工作；
  \item 用主机代码管理数据与调度，用设备核函数执行高并行计算；
  \item 用网格-线程块-线程层次和 $i=bT+t$ 建立执行到数据的映射；
  \item 用显式内存管理完成分配、传输、计算、回传与释放；
  \item 用向上取整启动和边界检查覆盖任意长度输入；
  \item 用 NVCC 编译主机与设备代码，并正确处理异步执行与错误状态；
  \item 先验证正确性，再进行性能优化。
\end{enumerate}

\begin{center}
\begin{tikzpicture}[
  node distance=7mm and 7mm,
  host/.style={draw=UIUCBlue,fill=SoftBlue,rounded corners,minimum height=9mm,
               text width=2.35cm,align=center,font=\small\sffamily},
  dev/.style={draw=UIUCOrange,fill=SoftOrange,rounded corners,minimum height=9mm,
              text width=2.35cm,align=center,font=\small\sffamily},
  flow/.style={-{Stealth[length=2.2mm]},thick,draw=Ink}
]
\node[host] (alloc) {分配设备内存\\allocate};
\node[host,right=of alloc] (h2d) {复制输入\\host to device};
\node[dev,right=of h2d] (kernel) {启动核函数\\kernel launch};
\node[host,right=of kernel] (d2h) {取回结果\\device to host};
\node[host,right=of d2h] (free) {释放资源\\free};
\draw[flow] (alloc) -- (h2d);
\draw[flow] (h2d) -- (kernel);
\draw[flow] (kernel) -- (d2h);
\draw[flow] (d2h) -- (free);
\node[below=7mm of kernel,dev,text width=5.4cm] (threads)
  {线程网格执行同一核函数；每个线程用 $i=bT+t$ 选择自己的数据。};
\draw[flow] (kernel) -- (threads);
\end{tikzpicture}
\end{center}

\begin{keybox}[title={关键公式}]
\[
i=\code{blockIdx.x}\cdot\code{blockDim.x}+\code{threadIdx.x},
\qquad
B=\frac{n+T-1}{T},
\]
\[
\code{if (i<n)}\ \text{保护边界},
\qquad
\code{<<<B,T>>>}\ \text{依次表示线程块数与每块线程数}.
\]
\end{keybox}
